\clearpage\SourceLeaf{434}\ArithmeticHead{427}{R}{KONVERGENZ DER ENTWICKLUNG DER MITTELPUNKTSGLEICHUNG.}
\begingroup\makeatletter\let\veqno\@@leqno\makeatother
\noindent{\small [Vierter Zettel]}\par
\begin{center}[5.]\end{center}
\noindent\hspace*{\parindent}Bestimmung des Zuges complexer Werthe von
\[E=p+qi;\]
\begin{equation}q=\log\operatorname{tg}(45^0+\tfrac12\theta),\tag*{[36]}\end{equation}
\noindent [für reelle Werte von \(M=E-f\sin E,\)]
\begin{equation}p-\frac{f\sin p}{\cos\theta}=M,\text{[\textsuperscript{*})]}\tag*{[37]}\end{equation}
\begin{equation}\frac{\log\operatorname{tg}(45^0+\tfrac12\theta)}{kf\,{.}\,\operatorname{tg}\theta}=\cos p,\tag*{[38]}\end{equation}
\noindent \(k\) [der] Modul[\textsuperscript{**})] d[er] \textsc{Brigg}[schen] Log[arithmen.]
\par\noindent\hspace*{\parindent}Für \(M=0;\quad p=0;\quad\theta=68.\ 27.\ 39{,}75\)
\[\begin{array}{r@{\;}l@{\qquad}l}
\log\operatorname{tg}\left(45^0+\dfrac\theta2\right)\;{.}\,{.}\,{.}&9{,}8577971\ 5&\text{Zahl}=0{,}7207708\\[4pt]
\operatorname{tg}\theta\;{.}\,{.}\,{.}&0{,}4037371\ 5&\\
kf\;{.}\,{.}\,{.}&9{,}4540600.&
\end{array}\]
\begin{center}\rule{.24\linewidth}{.4pt}\end{center}
\begingroup\setlength{\arraycolsep}{8pt}\renewcommand{\arraystretch}{1.2}
\[\begin{array}{c|r|l}
\theta&p&M\\\hline
71^0&-20^0\ 0'\ 0''&+\ 19.\ 25.\ 42\\
76&-37.\ 2.\ 0&+\ 56.\ 24.\ 12\text{[\textsuperscript{***})]}\\
81&-52.\ 4.\ 22&+137.\ 10.\ 24
\end{array}\]
\endgroup
\par\vfill\noindent\rule{.25\linewidth}{.4pt}\par\smallskip
{\footnotesize
\noindent\hspace*{\parindent}[\textsuperscript{*}) \textsc{Gauss} hat hier und im Folgenden die Exzentrizität mit \(e\) bezeichnet; wir haben in Übereinstimmung mit dem art. [2.] dafür \(f\) gesetzt.]\par
\noindent\hspace*{\parindent}[\textsuperscript{**}) Die Handschrift hat hier \textls[120]{Basis}.]\par
\noindent\hspace*{\parindent}[\textsuperscript{***}) In der Handschrift sind hier und in der Tabelle auf der folgenden Seite für \(\theta=76^0\) die unrichtigen Werte \(p=-35^0\,13'\,33''\), \(M-p=89^0\,29'\,3''\), \(M=54^0\,15'\,30''\) angegeben und errechnet. Nun ist der Logarithmus von \(\cos35^0\,13'\,33''\) gleich \(9{,}9\underline{1}216\), während der richtige Wert von \(\log_{10}\cos p\), wie er sich durch Subtraktion der Zahlen \(9{,}35622\) und \(9{,}45406\) (siehe auf der folgenden Seite) ergibt, \(9{,}9\underline{0}216\) lautet. \textsc{Gauss} hat also bei dieser Subtraktion ein Versehen begangen. Im Text sind überall die richtigen Zahlen eingesetzt.]\par}
\noindent\hfill{\small 54\textsuperscript{*}}\hspace*{.1\linewidth}\par
\clearpage\SourceLeaf{435}\ArithmeticHead{428}{L}{ANALYSIS. NACHLASS.}
\begingroup\setlength{\arraycolsep}{8pt}\renewcommand{\arraystretch}{1.3}
\[\begin{array}{r@{\;}c|c}
&\theta=76^0&\theta=81^0\\\cline{2-3}
\log\operatorname{tg}\left(45^0+\dfrac\theta2\right)&9{,}95945&0{,}04298\\[3pt]
\operatorname{tg}\theta&0{,}60323&0{,}80029\\
&9{,}35622&9{,}24269\\
kf&\underline{9{,}45406}&\underline{9{,}45406}\\
{}[\cos p&9{,}90216&9{,}78863]\\
p&-37.\ 2.\ 0\text{[\textsuperscript{*})]}&[-52.\ 4.\ 22]\\
\sin p&9{,}77980&9{,}89696\\
f\text{ in Sec.}&\underline{5{,}13070}&\underline{5{,}13070}\\
&4{,}91050&5{,}02766\\
\cos\theta&9{,}38368&9{,}19433\\
{}[M-p\text{ in Sec. }=]&336372&681286\\
&93.\ 26.\ 12\text{[\textsuperscript{*})]}&189.\ 14.\ 46
\end{array}\]
\endgroup
\par\noindent\rule{.25\linewidth}{.4pt}\par\smallskip
{\footnotesize\noindent\hspace*{\parindent}[\textsuperscript{*}) Siehe die Fußnote \textsuperscript{***}) auf der vorigen Seite.]\par}
\vspace{2cm}\begin{center}\rule{.32\linewidth}{.4pt}\end{center}
\clearpage\endgroup
